\documentclass[11pt,a4paper]{article}
\usepackage[utf8]{inputenc}
\usepackage[margin=2.4cm]{geometry}
\usepackage{amsmath,amssymb}
\usepackage{booktabs}
\usepackage{xcolor}
\usepackage{graphicx}
\usepackage{listings}
\usepackage{enumitem}
\usepackage{microtype}
\usepackage[colorlinks=true,linkcolor=blue!50!black,urlcolor=blue!50!black]{hyperref}

\definecolor{codegray}{gray}{0.45}
\definecolor{codeblue}{rgb}{0.1,0.2,0.6}
\lstset{
  language=Python,
  basicstyle=\ttfamily\footnotesize,
  keywordstyle=\color{codeblue},
  commentstyle=\color{codegray}\itshape,
  stringstyle=\color{teal},
  showstringspaces=false,
  breaklines=true,
  frame=single,
  rulecolor=\color{black!25},
  numbers=left,
  numberstyle=\tiny\color{codegray},
  literate={θ}{$\theta$}1 {ξ}{$\xi$}1,
}

\newcommand{\session}[3]{%
  \title{\textbf{#1}\\[2pt]\large #2}
  \author{astro-lab \,--\, #3}
  \date{}
  \maketitle
  \vspace{-1.5em}
}
\newcommand{\todo}[1]{\textcolor{red!60!black}{\textbf{[#1]}}}
\usepackage{framed}
\newenvironment{keyidea}{%
  \def\FrameCommand{\color{blue!40!black}\vrule width 2.5pt\hspace{8pt}\color{black}}%
  \MakeFramed{\advance\hsize-\width\FrameRestore}}{\endMakeFramed}

\newcommand{\dd}{\mathrm{d}}
\newcommand{\pc}{\,\mathrm{pc}}
\newcommand{\Msun}{M_\odot}
\newcommand{\Rsun}{R_\odot}
\setlist[itemize]{itemsep=2pt,topsep=4pt}
\setlist[enumerate]{itemsep=3pt,topsep=4pt}
\setlength{\parindent}{0pt}
\setlength{\parskip}{6pt}

\begin{document}
\session{Session 6: Computational Hydrodynamics in 1D}{Euler equations, finite volume, the HLL flux, the Sod shock tube}{reading notes}

\textbf{Status:} reading notes and exercise specification. The code for this session (\texttt{06\_hydro\_1d/}) is not in the repository yet; the exercises below define what to build.

\section*{Goals}
Write a finite-volume solver for the 1D Euler equations, validate it on the Sod shock tube against the exact solution, and measure how its error behaves as the grid is refined.

\section{The Euler equations}
Gas dynamics without viscosity, conduction or gravity is a system of conservation laws,
\[
\frac{\partial\mathbf U}{\partial t}+\frac{\partial\mathbf F(\mathbf U)}{\partial x}=0,
\qquad
\mathbf U=\begin{pmatrix}\rho\\ \rho u\\ E\end{pmatrix},\quad
\mathbf F=\begin{pmatrix}\rho u\\ \rho u^2+P\\ u(E+P)\end{pmatrix},
\]
with total energy density $E=\dfrac{P}{\gamma-1}+\tfrac12\rho u^2$ for an ideal gas, and sound speed $c=\sqrt{\gamma P/\rho}$. Because the equations are in conservation form, mass, momentum and energy are conserved to round-off by any scheme that updates cells by flux differences. That is the main reason to use finite volume rather than finite differences on primitive variables.

\section{Finite-volume update}
Divide the domain into cells of width $\Delta x$ with cell averages $\mathbf U_i$. Then
\[
\mathbf U_i^{n+1}=\mathbf U_i^n-\frac{\Delta t}{\Delta x}\left(\mathbf F_{i+1/2}-\mathbf F_{i-1/2}\right),
\]
where $\mathbf F_{i+1/2}$ is a numerical flux computed from the states on the two sides of the interface. Stability requires the CFL condition
\[
\Delta t=C\,\frac{\Delta x}{\max_i\left(|u_i|+c_i\right)},\qquad C\le1 \ (\text{use }0.4\text{--}0.8).
\]

\subsection*{The HLL flux}
Estimate the fastest left- and right-moving wave speeds at an interface (Davis estimates):
\[
S_L=\min(u_L-c_L,\ u_R-c_R),\qquad S_R=\max(u_L+c_L,\ u_R+c_R).
\]
Then
\[
\mathbf F^{\mathrm{HLL}}=
\begin{cases}
\mathbf F_L & S_L\ge0,\\[3pt]
\dfrac{S_R\mathbf F_L-S_L\mathbf F_R+S_LS_R(\mathbf U_R-\mathbf U_L)}{S_R-S_L} & S_L<0<S_R,\\[10pt]
\mathbf F_R & S_R\le0.
\end{cases}
\]
HLL is robust and simple but smears contact discontinuities because it ignores the middle wave. HLLC restores it. Lax--Friedrichs is the simplest flux and the most diffusive; it is a good first implementation to get the bookkeeping right.

\begin{keyidea}
\textbf{First-order versus second-order.} Using the cell values directly on both sides of an interface gives a first-order scheme: stable but diffusive. Reconstructing a linear profile in each cell with a slope limiter (minmod, MC) and a half-step predictor gives second order in smooth regions without spurious oscillations at shocks.
\end{keyidea}

\section{The Sod shock tube}
Initial conditions on $x\in[0,1]$, diaphragm at $x=0.5$, $\gamma=1.4$:
\[
(\rho,u,P)=(1,\,0,\,1)\ \ (x<0.5),\qquad (\rho,u,P)=(0.125,\,0,\,0.1)\ \ (x>0.5).
\]
At $t=0.2$ the exact solution has a left-moving rarefaction fan, a contact discontinuity and a right-moving shock. Standard values for the intermediate states:
\begin{center}
\begin{tabular}{@{}lr@{}}
\toprule
Quantity & Value \\ \midrule
$P_*$ & 0.30313 \\
$u_*$ & 0.92745 \\
$\rho_{*L}$ (left of contact) & 0.42632 \\
$\rho_{*R}$ (right of contact) & 0.26557 \\
\bottomrule
\end{tabular}
\end{center}
At $t=0.2$ the features sit at approximately $x\approx0.263$ (rarefaction head), $0.486$ (rarefaction tail), $0.686$ (contact) and $0.850$ (shock).

\section{Implementation plan}
\begin{enumerate}
\item Arrays of conserved variables with two ghost cells; transmissive boundaries.
\item Functions \texttt{prim\_to\_cons}, \texttt{cons\_to\_prim}, \texttt{flux}, \texttt{hll\_flux}.
\item Time loop with the CFL timestep, stopping exactly at $t=0.2$.
\item Check: total mass, momentum and energy are conserved to round-off until waves reach the boundary.
\item Compare density, velocity and pressure with the exact solution (write a Riemann solver or use a published reference solution).
\item Convergence study: $N=100,200,400,800,1600$; plot the $L_1$ density error versus $\Delta x$ on log--log axes.
\end{enumerate}

\section{Exercises}
\begin{enumerate}
\item Implement Lax--Friedrichs first and compare against HLL; describe the difference at the contact.
\item Add MUSCL reconstruction with a minmod limiter and measure the improvement.
\item Explain why the measured $L_1$ convergence rate at a solution with discontinuities is lower than the scheme's formal order.
\item Test a case with a strong shock-tube such as the Lax problem or a double rarefaction (123 problem) and note any failures (negative pressure) and how to guard against them.
\item Connect to Session 3: add a static gravitational source term and check that an isothermal atmosphere stays in hydrostatic balance.
\end{enumerate}

\section*{References}
\begin{itemize}
\item Toro, \textit{Riemann Solvers and Numerical Methods for Fluid Dynamics}, ch.\ 4, 10
\item LeVeque, \textit{Finite Volume Methods for Hyperbolic Problems}
\item Zingale, \textit{Computational Hydrodynamics for Astrophysics} (lecture notes; search for the title)
\item Sod (1978), \textit{J.\ Comput.\ Phys.} 27, 1
\end{itemize}
\end{document}
